\chapter{核心查询引擎 QueryEngine}

\section{QueryEngine 概述}

\code{QueryEngine} 是 Claude Code 的核心类，定义在 \code{src/QueryEngine.ts} 中（约 46,000 行）。它拥有整个对话生命周期的状态，管理与 Anthropic API 的通信、工具调用循环、流式响应处理等。

\subsection{设计定位}

\begin{lstlisting}[language=TypeScript,breaklines=true]
/**
 * QueryEngine 拥有查询生命周期和会话状态。
 * 它将核心逻辑从 ask() 提取为独立类，既可用于 headless/SDK 路径，
 * 也可（在未来阶段）用于 REPL。
 *
 * 每个对话一个 QueryEngine。每次 submitMessage() 调用在同一对话中
 * 启动一个新的回合（turn）。状态（消息、文件缓存、使用量等）
 * 在回合之间持久存在。
 */
export class QueryEngine {
  private config: QueryEngineConfig
  private mutableMessages: Message[]
  private abortController: AbortController
  private permissionDenials: SDKPermissionDenial[]
  private totalUsage: NonNullableUsage
  private hasHandledOrphanedPermission = false
  private readFileState: FileStateCache
  private discoveredSkillNames = new Set<string>()
  private loadedNestedMemoryPaths = new Set<string>()

  constructor(config: QueryEngineConfig) {
    this.config = config
    this.mutableMessages = config.initialMessages ?? []
    this.abortController = config.abortController ?? createAbortController()
    this.permissionDenials = []
    this.readFileState = config.readFileCache
    this.totalUsage = EMPTY_USAGE
  }
}
\end{lstlisting}

\subsection{配置接口}

\begin{lstlisting}[language=TypeScript,breaklines=true]
export type QueryEngineConfig = {
  cwd: string                           // 工作目录
  tools: Tools                          // 可用工具列表
  commands: Command[]                   // 可用命令列表
  mcpClients: MCPServerConnection[]     // MCP 服务器连接
  agents: AgentDefinition[]             // 代理定义
  canUseTool: CanUseToolFn              // 工具权限检查函数
  getAppState: () => AppState           // 全局状态获取
  setAppState: (f) => void              // 全局状态更新
  initialMessages?: Message[]           // 初始消息历史
  readFileCache: FileStateCache         // 文件状态缓存
  customSystemPrompt?: string           // 自定义系统提示
  appendSystemPrompt?: string           // 追加系统提示
  userSpecifiedModel?: string           // 用户指定模型
  fallbackModel?: string                // 备用模型
  thinkingConfig?: ThinkingConfig       // 思维模式配置
  maxTurns?: number                     // 最大回合数
  maxBudgetUsd?: number                 // 最大预算（美元）
  taskBudget?: { total: number }        // 任务预算
  jsonSchema?: Record<string, unknown>  // 结构化输出 Schema
  verbose?: boolean                     // 详细模式
  replayUserMessages?: boolean          // 重放用户消息
  handleElicitation?: (...)  => Promise  // URL 引出处理
  includePartialMessages?: boolean      // 包含部分消息
  setSDKStatus?: (status) => void       // SDK 状态回调
  abortController?: AbortController     // 中止控制器
  orphanedPermission?: OrphanedPermission // 孤立权限
  snipReplay?: (...)  => ... | undefined  // Snip 重放处理器
}
\end{lstlisting}

\section{submitMessage：核心方法}

\code{submitMessage} 是一个 \textbf{AsyncGenerator}——它通过 \code{yield} 逐步产出流式事件，使调用方能实时获取 API 响应的每一部分。

\begin{lstlisting}[language=TypeScript,breaklines=true]
async *submitMessage(
  prompt: string | ContentBlockParam[],
  options?: { uuid?: string; isMeta?: boolean },
): AsyncGenerator<SDKMessage, void, unknown> {
\end{lstlisting}

\subsection{执行流程}

\begin{lstlisting}[language=TypeScript,breaklines=true]
submitMessage(prompt)
    │
    ▼
1. 清理状态 & 设置工作目录
    │
    ▼
2. 包装 canUseTool（追踪权限拒绝）
    │
    ▼
3. 解析模型 & 思维配置
    │
    ▼
4. 获取 System Prompt
    ├── fetchSystemPromptParts()
    ├── 默认系统提示 (defaultSystemPrompt)
    ├── 用户上下文 (baseUserContext)
    └── 系统上下文 (systemContext)
    │
    ▼
5. 构建 ProcessUserInputContext
    │
    ▼
6. 处理用户输入 → processUserInput()
    ├── 斜杠命令检测
    ├── 附件处理
    └── 生成用户消息
    │
    ▼
7. 调用查询循环 → query()
    ├── 构建完整 System Prompt
    ├── 调用 Anthropic API（流式）
    ├── 处理 Tool Use 请求
    ├── 执行工具 → 返回结果
    ├── 自动压缩（如果需要）
    └── 循环直到模型停止
    │
    ▼
8. 产出最终状态 & 使用量统计
\end{lstlisting}

\subsection{权限追踪包装}

\begin{lstlisting}[language=TypeScript,breaklines=true]
const wrappedCanUseTool: CanUseToolFn = async (
  tool, input, toolUseContext, assistantMessage, toolUseID, forceDecision,
) => {
  const result = await canUseTool(
    tool, input, toolUseContext, assistantMessage, toolUseID, forceDecision,
  )

  // 追踪拒绝以供 SDK 报告
  if (result.behavior !== 'allow') {
    this.permissionDenials.push({
      tool_name: sdkCompatToolName(tool.name),
      tool_use_id: toolUseID,
      tool_input: input,
    })
  }

  return result
}
\end{lstlisting}

每次工具调用的权限决定都被包装和追踪。被拒绝的工具调用会被记录，最终通过 SDK 事件报告给调用方。

\subsection{思维模式配置}

\begin{lstlisting}[language=TypeScript,breaklines=true]
const initialThinkingConfig: ThinkingConfig = thinkingConfig
  ? thinkingConfig
  : shouldEnableThinkingByDefault() !== false
    ? { type: 'adaptive' }   // 自适应思维
    : { type: 'disabled' }   // 禁用思维
\end{lstlisting}

Claude Code 支持三种思维模式：
\begin{itemize}[nosep]
  \item \textbf{\code{adaptive}}：模型自行决定是否需要深度思考
  \item \textbf{\code{enabled}}：始终启用扩展思考
  \item \textbf{\code{disabled}}：禁用扩展思考
\end{itemize}

\section{查询循环 (query.ts)}

\code{query.ts} 实现了与 API 交互的核心循环——这是 Claude Code 最复杂的控制流之一。

\subsection{循环状态}

\begin{lstlisting}[language=TypeScript,breaklines=true]
type State = {
  messages: Message[]                              // 消息历史
  toolUseContext: ToolUseContext                    // 工具使用上下文
  autoCompactTracking: AutoCompactTrackingState    // 自动压缩追踪
  maxOutputTokensRecoveryCount: number             // 输出 Token 恢复计数
  hasAttemptedReactiveCompact: boolean             // 是否尝试过响应式压缩
  maxOutputTokensOverride: number | undefined      // 输出 Token 覆盖值
  pendingToolUseSummary: Promise<...> | undefined  // 待处理工具使用摘要
  stopHookActive: boolean | undefined              // 停止钩子是否激活
  turnCount: number                                // 回合计数
  transition: Continue | undefined                 // 上一次迭代的继续原因
}
\end{lstlisting}

\subsection{循环的核心逻辑}

\begin{lstlisting}[language=TypeScript,breaklines=true]
export async function* query(params: QueryParams): AsyncGenerator<...> {
  const consumedCommandUuids: string[] = []
  const terminal = yield* queryLoop(params, consumedCommandUuids)
  
  for (const uuid of consumedCommandUuids) {
    notifyCommandLifecycle(uuid, 'completed')
  }
  return terminal
}

async function* queryLoop(params, consumedCommandUuids) {
  // ... 初始化状态 ...
  
  while (true) {
    // 1. 技能发现预取
    const pendingSkillPrefetch = skillPrefetch?.startSkillDiscoveryPrefetch(...)
    
    // 2. 发出请求开始事件
    yield { type: 'stream_request_start' }
    
    // 3. 获取压缩边界后的消息
    let messagesForQuery = [...getMessagesAfterCompactBoundary(messages)]
    
    // 4. 应用工具结果预算
    messagesForQuery = await applyToolResultBudget(messagesForQuery, ...)
    
    // 5. Snip 压缩（如果启用）
    if (feature('HISTORY_SNIP')) {
      const snipResult = snipModule!.snipCompactIfNeeded(messagesForQuery)
      messagesForQuery = snipResult.messages
    }
    
    // 6. 微压缩
    const microcompactResult = await deps.microcompact(messagesForQuery, ...)
    messagesForQuery = microcompactResult.messages
    
    // 7. 上下文折叠（如果启用）
    if (feature('CONTEXT_COLLAPSE') && contextCollapse) {
      const collapseResult = await contextCollapse.applyCollapsesIfNeeded(...)
      messagesForQuery = collapseResult.messages
    }
    
    // 8. 构建完整系统提示
    const fullSystemPrompt = asSystemPrompt(
      appendSystemContext(systemPrompt, systemContext)
    )
    
    // 9. 调用 API 并处理响应
    // ... 流式处理 + 工具调用 ...
    
    // 10. 判断是否继续循环
    // - 如果有工具调用 → 继续
    // - 如果模型结束 → 退出
    // - 如果达到 max_output_tokens → 恢复尝试
  }
}
\end{lstlisting}

\subsection{多层压缩管线}

查询循环中的消息在发送到 API 之前，会经过一个\textbf{多层压缩管线}：

\begin{lstlisting}[language=TypeScript,breaklines=true]
原始消息历史
    │
    ▼
① 工具结果预算 (applyToolResultBudget)
    │  → 替换超大工具结果为摘要
    ▼
② Snip 压缩 (snipCompactIfNeeded)
    │  → 移除旧的对话片段
    ▼
③ 微压缩 (microcompact)
    │  → 压缩单个工具输出的冗余内容
    ▼
④ 上下文折叠 (applyCollapsesIfNeeded)
    │  → 将多轮交互折叠为摘要
    ▼
⑤ 自动压缩 (autoCompact)
    │  → 当接近上下文窗口限制时触发
    ▼
发送到 API
\end{lstlisting}

\subsection{max\_output\_tokens 恢复机制}

当模型输出被截断时，Claude Code 会自动尝试恢复：

\begin{lstlisting}[language=TypeScript,breaklines=true]
const MAX_OUTPUT_TOKENS_RECOVERY_LIMIT = 3

function isWithheldMaxOutputTokens(msg: Message | undefined): boolean {
  return msg?.type === 'assistant' && msg.apiError === 'max_output_tokens'
}
\end{lstlisting}

最多尝试 3 次恢复。在恢复过程中，SDK 调用方不会收到错误消息——错误被"扣留"（withheld），直到确认恢复循环是否成功。

\section{思维规则}

源码中包含一段幽默的注释，描述了处理思维块（thinking blocks）的规则：

\begin{lstlisting}[language=TypeScript,breaklines=true]
/**
 * 思维规则冗长而繁琐。它们需要大量的长时间思考和深度冥想，
 * 才能让一个巫师将其理解透彻。
 *
 * 规则如下：
 * 1. 包含 thinking 或 redacted_thinking 块的消息必须属于
 *    max_thinking_length > 0 的查询
 * 2. thinking 块不能是块中的最后一条消息
 * 3. 思维块必须在助手轨迹期间保持完整
 *    （单个回合，或如果该回合包含 tool_use 块，
 *     则也包括其后续的 tool_result 和跟随的助手消息）
 *
 * 好好记住这些规则，年轻的巫师。因为它们是思维的规则，
 * 而思维的规则就是宇宙的规则。如果你不遵守这些规则，
 * 你将受到一整天调试和拔头发的惩罚。
 */
\end{lstlisting}

这段注释反映了 Anthropic 工程师在调试思维块处理逻辑时遇到的痛苦经验。

\section{ToolUseContext：工具执行上下文}

\code{ToolUseContext} 是传递给每个工具调用的丰富上下文对象：

\begin{lstlisting}[language=TypeScript,breaklines=true]
export type ToolUseContext = {
  options: {
    commands: Command[]
    debug: boolean
    mainLoopModel: string
    tools: Tools
    verbose: boolean
    thinkingConfig: ThinkingConfig
    mcpClients: MCPServerConnection[]
    mcpResources: Record<string, ServerResource[]>
    isNonInteractiveSession: boolean
    agentDefinitions: AgentDefinitionsResult
    maxBudgetUsd?: number
    customSystemPrompt?: string
    appendSystemPrompt?: string
    refreshTools?: () => Tools
  }
  abortController: AbortController
  readFileState: FileStateCache
  getAppState(): AppState
  setAppState(f: (prev: AppState) => AppState): void
  setAppStateForTasks?: (...)  => void
  handleElicitation?: (...)  => Promise<ElicitResult>
  setToolJSX?: SetToolJSXFn
  addNotification?: (notif: Notification) => void
  appendSystemMessage?: (msg) => void
  sendOSNotification?: (opts) => void
  messages: Message[]
  fileReadingLimits?: { maxTokens?: number; maxSizeBytes?: number }
  globLimits?: { maxResults?: number }
  toolDecisions?: Map<string, { source; decision; timestamp }>
  queryTracking?: QueryChainTracking
  contentReplacementState?: ContentReplacementState
  renderedSystemPrompt?: SystemPrompt
  // ... 更多字段
}
\end{lstlisting}

关键设计点：

\begin{enumerate}[nosep]
  \item \textbf{\code{readFileState}}：文件状态 LRU 缓存。避免重复读取同一文件。
  \item \textbf{\code{setAppStateForTasks}}：允许嵌套子代理更新根状态存储，即使子代理的 \code{setAppState} 是空操作。
  \item \textbf{\code{contentReplacementState}}：工具结果预算的替换状态，跨轮次持久化。
  \item \textbf{\code{renderedSystemPrompt}}：父代理的渲染系统提示的冻结快照，用于 fork 子代理共享缓存。
\end{enumerate}

\section{Token 预算追踪}

\begin{lstlisting}[language=TypeScript,breaklines=true]
const budgetTracker = feature('TOKEN_BUDGET') 
  ? createBudgetTracker() 
  : null
\end{lstlisting}

当 \code{TOKEN\_BUDGET} Feature Flag 启用时，查询循环会追踪 Token 消耗并与预算比较。

\textbf{任务预算 (Task Budget)}：通过 API 的 \code{output\_config.task\_budget} 参数设置整个代理回合的预算。每次压缩后需要重新计算剩余预算。

\section{设计启示}

\subsection{AsyncGenerator 模式}

使用 \code{async function*} 实现流式数据处理是一个优雅的设计。调用方可以通过 \code{for await...of} 消费流，也可以使用 \code{.next()} 精确控制步进。

\subsection{多层压缩管线}

五层压缩看似复杂，但每层解决的问题不同：
\begin{itemize}[nosep]
  \item \textbf{工具结果预算}：控制单个工具输出的大小
  \item \textbf{Snip}：移除老旧的历史
  \item \textbf{微压缩}：局部优化
  \item \textbf{上下文折叠}：语义级别的压缩
  \item \textbf{自动压缩}：全局兜底
\end{itemize}

\subsection{恢复机制}

对 \code{max\_output\_tokens} 的自动恢复展示了一个重要原则：\textbf{AI 应用应该对 API 的边界条件有鲁棒的处理策略}，而不是简单地向用户抛出错误。

\bigskip\noindent\rule{\textwidth}{0.4pt}\bigskip

\textbf{下一章}：\href{./chapter-05-context.md}{上下文处理与 System Prompt}
